% GENERATED FILE. Edit canonical lessons, then run npm run generate:books.
% canonical-source-hash: fnv1a64:17c0c46172d7f8b1
% canonical-lessons: TE-C163-anumanincu, TE-C163-grahincu, TE-C163-pattincuko, TE-C163-potipadu, TE-C163-tannu

\chapter{To suspect, To realize, To pay attention to, To compete, To kick}
\label{ch:te-to-suspect-to-realize-to-pay-attention-to-to-compete-to-kick}
\hlchaptermodality{163}

\begin{chapteropening}
\textbf{By the end of this chapter:} \emph{I can say to suspect, to realize, to pay attention to, to compete and to kick.}

Complete the last lesson of chapter 163: I can say to suspect, to realize, to pay attention to, to compete and to kick.
\end{chapteropening}

\section[\texorpdfstring{\te{అనుమానించు} (anumāniñcu)}{అనుమానించు (anumāniñcu)}]{\te{అనుమానించు} (anumāniñcu) — to suspect, to doubt}

\label{lesson:TE-C163-anumanincu}



\begin{quote}
Before the new one: say the Telugu for to praise, then the Telugu for to describe.
\end{quote}

\subsection*{You'll want to know: \te{అనుమానించు}}
\textbf{\te{అనుమానించు}} — \emph{anumāniñcu} — \textquotedblleft{}to suspect, to doubt\textquotedblright{}.

\begin{grammarlens}[title={What you've built}]
One more word for this chapter.
\end{grammarlens}

\subsection*{Your turn}
\begin{itemize}
\raggedright
  \item \emph{Say it:} \emph{anumāniñcu}
  \item \emph{Say it:} \emph{anumāniñcu}, once more
  \item \emph{Say it:} say \emph{anumāniñcu}
  \item \emph{Recall:} say the Telugu for to praise, then the Telugu for to describe, then say \emph{anumāniñcu} again
\end{itemize}

\begin{tcolorbox}[breakable,colback=teal!4,colframe=teal!35!black,title={Before you move on}]
What does \te{అనుమానించు} mean? (\textquotedblleft{}To suspect, to doubt\textquotedblright{}.) Say it once more.
\end{tcolorbox}
\section[\texorpdfstring{\te{గ్రహించు} (grahiñcu)}{గ్రహించు (grahiñcu)}]{\te{గ్రహించు} (grahiñcu) — to realize, to grasp}

\label{lesson:TE-C163-grahincu}



\begin{quote}
Before the new one: say the Telugu for to describe, then the Telugu for to suspect.
\end{quote}

\subsection*{You'll want to know: \te{గ్రహించు}}
\textbf{\te{గ్రహించు}} — \emph{grahiñcu} — \textquotedblleft{}to realize, to grasp\textquotedblright{}.

\begin{grammarlens}[title={What you've built}]
One more word for this chapter.
\end{grammarlens}

\subsection*{Your turn}
\begin{itemize}
\raggedright
  \item \emph{Say it:} \emph{grahiñcu}
  \item \emph{Say it:} \emph{grahiñcu}, once more
  \item \emph{Say it:} say \emph{grahiñcu}
  \item \emph{Recall:} say the Telugu for to describe, then the Telugu for to suspect, then say \emph{grahiñcu} again
\end{itemize}

\begin{tcolorbox}[breakable,colback=teal!4,colframe=teal!35!black,title={Before you move on}]
What does \te{గ్రహించు} mean? (\textquotedblleft{}To realize, to grasp\textquotedblright{}.) Say it once more.
\end{tcolorbox}
\section[\texorpdfstring{\te{పట్టించుకో} (paṭṭiñcukō)}{పట్టించుకో (paṭṭiñcukō)}]{\te{పట్టించుకో} (paṭṭiñcukō) — to pay attention to, to care about}

\label{lesson:TE-C163-pattincuko}



\begin{quote}
Before the new one: say the Telugu for to suspect, then the Telugu for to realize.
\end{quote}

\subsection*{You'll want to know: \te{పట్టించుకో}}
\textbf{\te{పట్టించుకో}} — \emph{paṭṭiñcukō} — \textquotedblleft{}to pay attention to, to care about\textquotedblright{}.

\begin{grammarlens}[title={What you've built}]
One more word for this chapter.
\end{grammarlens}

\subsection*{Your turn}
\begin{itemize}
\raggedright
  \item \emph{Say it:} \emph{paṭṭiñcukō}
  \item \emph{Say it:} \emph{paṭṭiñcukō}, once more
  \item \emph{Say it:} say \emph{paṭṭiñcukō}
  \item \emph{Recall:} say the Telugu for to suspect, then the Telugu for to realize, then say \emph{paṭṭiñcukō} again
\end{itemize}

\begin{tcolorbox}[breakable,colback=teal!4,colframe=teal!35!black,title={Before you move on}]
What does \te{పట్టించుకో} mean? (\textquotedblleft{}To pay attention to, to care about\textquotedblright{}.) Say it once more.
\end{tcolorbox}
\section[\texorpdfstring{\te{పోటీపడు} (pōṭīpaḍu)}{పోటీపడు (pōṭīpaḍu)}]{\te{పోటీపడు} (pōṭīpaḍu) — to compete}

\label{lesson:TE-C163-potipadu}



\begin{quote}
Before the new one: say the Telugu for to realize, then the Telugu for to pay attention to.
\end{quote}

\subsection*{You'll want to know: \te{పోటీపడు}}
\textbf{\te{పోటీపడు}} — \emph{pōṭīpaḍu} — \textquotedblleft{}to compete\textquotedblright{}.

\begin{grammarlens}[title={What you've built}]
One more word for this chapter.
\end{grammarlens}

\subsection*{Your turn}
\begin{itemize}
\raggedright
  \item \emph{Say it:} \emph{pōṭīpaḍu}
  \item \emph{Say it:} \emph{pōṭīpaḍu}, once more
  \item \emph{Say it:} say \emph{pōṭīpaḍu}
  \item \emph{Recall:} say the Telugu for to realize, then the Telugu for to pay attention to, then say \emph{pōṭīpaḍu} again
\end{itemize}

\begin{tcolorbox}[breakable,colback=teal!4,colframe=teal!35!black,title={Before you move on}]
What does \te{పోటీపడు} mean? (\textquotedblleft{}To compete\textquotedblright{}.) Say it once more.
\end{tcolorbox}
\section[\texorpdfstring{\te{తన్ను} (tannu)}{తన్ను (tannu)}]{\te{తన్ను} (tannu) — to kick}

\label{lesson:TE-C163-tannu}



\begin{quote}
Before the new one: say the Telugu for to pay attention to, then the Telugu for to compete.
\end{quote}

\subsection*{You'll want to know: \te{తన్ను}}
\textbf{\te{తన్ను}} — \emph{tannu} — \textquotedblleft{}to kick\textquotedblright{}.

\begin{grammarlens}[title={What you've built}]
One more word for this chapter.
\end{grammarlens}

\subsection*{Your turn}
\begin{itemize}
\raggedright
  \item \emph{Say it:} \emph{tannu}
  \item \emph{Say it:} \emph{tannu}, once more
  \item \emph{Say it:} say \emph{tannu}
  \item \emph{Recall:} say the Telugu for to pay attention to, then the Telugu for to compete, then say \emph{tannu} again
\end{itemize}

\begin{tcolorbox}[breakable,colback=teal!4,colframe=teal!35!black,title={Before you move on}]
What does \te{తన్ను} mean? (\textquotedblleft{}To kick\textquotedblright{}.) Say it once more.
\end{tcolorbox}
